% An IEEEtran conference paper without line numbers, with a listing numbered in the margin and a table that numbers
% its rows.
\pdftrailerid{}
\pdfinfoomitdate=1
\pdfsuppressptexinfo=-1
\documentclass[10pt,conference]{IEEEtran}
\usepackage{listings,lipsum}
\lstset{basicstyle=\ttfamily\scriptsize, numbers=left, numberstyle=\tiny, xleftmargin=1em}
\begin{document}
\title{Ranking Flaky Test Failures with Build Histories}
\author{\IEEEauthorblockN{Anonymous Author(s)}}
\maketitle
\begin{abstract}
Flaky tests pass and fail on the same code. The tool of this paper ranks the failures of a build by the chance that
each failure is flaky.
\end{abstract}
\section{Introduction}
\lipsum[1-2]
\begin{lstlisting}[language=Java]
@Test
public void testCacheOne() {
    assertEquals(1, cache.size());
}
@Test
public void testCacheTwo() {
    assertEquals(2, cache.size());
}
@Test
public void testCacheThree() {
    assertEquals(3, cache.size());
}
@Test
public void testCacheFour() {
    assertEquals(4, cache.size());
}
@Test
public void testCacheFive() {
    assertEquals(5, cache.size());
}
\end{lstlisting}
\begin{table}[t]
\caption{Projects of the study}
\centering
\begin{tabular}{rl}
1 & alpha \\ 2 & beta \\ 3 & gamma \\ 4 & delta \\ 5 & epsilon \\ 6 & zeta \\ 7 & eta \\ 8 & theta \\
9 & iota \\ 10 & kappa \\ 11 & lambda \\ 12 & mu \\ 13 & nu \\ 14 & xi \\ 15 & omicron \\ 16 & pi \\
\end{tabular}
\end{table}
\lipsum[3-7]
\end{document}
